\section{Related Work}\label{sec:related}

\textbf{Scheduling.} Scheduling under constraints is a core engineering problem, from job shops and production lines to staff rostering~\cite{pinedo2022}. University timetabling is usually studied from the institution's side, assigning rooms and times so that curricula do not clash~\cite{alghamdi2020}, with recent work adding individual study plans~\cite{steiner2025} and layered preferences~\cite{dunke2023}. On the student's side, planners already exist: UNSW's Notangles, for example, offers automatic timetabling from a unit list~\cite{notangles}. Among the tools and studies reviewed here, none takes the student's situation as free text, asks back when unsure, and explains its choices; those are what UniPlan adds, together with a ranking learned from preference comparisons.

\textbf{From text to constraints.} LLMs are widely used to turn free text into structured fields~\cite{xu2024ie}. Closest to our parser, the NL4Opt competition~\cite{ramamonjison2023nl4opt} and OptiMUS~\cite{ahmaditeshnizi2024optimus} turn written descriptions of optimisation problems into models a solver can run. We do the same for a narrower problem, with a fixed set of conditions that the user confirms before the solver runs. LoRA (low-rank adaptation)~\cite{hu2022lora} and its 4-bit quantised form QLoRA~\cite{dettmers2023qlora} let us fine-tune a 7B open model such as Qwen2.5~\cite{qwen2025} on one consumer GPU.

\textbf{Learning preferences.} Our ranker is the linear case, with student terms, of pairwise preference models going back to Bradley and Terry~\cite{bradley1952} and RankNet~\cite{burges2005ranknet}. LLMs can rival crowd workers as annotators~\cite{gilardi2023} and agree with human judges about as often as humans agree with each other, but they favour whichever option is shown first~\cite{zheng2023judge}, so we ask every comparison in both orders.
